%----------
%   IMPORTANTE
%----------

% Esta plantilla está basada en las recomendaciones de la guía "Trabajo fin de Grado: Escribir el TFG", que encontrarás en https://uc3m.libguides.com/tesis/escribir
% Esta Guía contiene recomendaciones de la Biblioteca basadas principalmente en estilos APA e IEEE, pero debes seguir siempre las orientaciones de tu Director de tesis y la normativa de la Escuela de Doctorado.

%----------
%	CONFIGURACIÓN DEL DOCUMENTO
%----------

\documentclass[12pt]{report} % fuente a 12pt

% MÁRGENES: 2,5 cm sup. e inf.; 3 cm izdo. y dcho.
\usepackage[
a4paper,
vmargin=2.5cm,
hmargin=3cm
]{geometry}

% INTERLINEADO: Estrecho (6 ptos./interlineado 1,15) o Moderado (6 ptos./interlineado 1,5)
\renewcommand{\baselinestretch}{1.15}
\parskip=6pt

% Soporte para GENERAR PDF/A --es importante de cara a su inclusión en e-Archivo porque es el formato óptimo de preservación y a la generación de metadatos, tal y como se describe en http://uc3m.libguides.com/ld.php?content_id=31389625. 

% En la plantilla incluimos el archivo OUTPUT.XMPDATA. Puedes descargar este archivo e incluir los metadatos que se incorporarán al archivo PDF cuando compiles el archivo memoria.tex. Después vuelve a subirlo a tu proyecto. 
\usepackage[a-1b]{pdfx}

% ENLACES
\usepackage{hyperref}
\hypersetup{colorlinks=true,
	linkcolor=black, % enlaces a partes del documento (p.e. índice) en color negro
	urlcolor=blue} % enlaces a recursos fuera del documento en azul

% EXPRESIONES MATEMÁTICAS
\usepackage{amsmath,amssymb,amsfonts,amsthm}

% Codificación caracteres
\usepackage{txfonts} 
\usepackage[T1]{fontenc}
\usepackage[utf8]{inputenc}

% Definición idioma español
\usepackage[spanish, es-tabla]{babel} 
\usepackage[babel, spanish=spanish]{csquotes}
\AtBeginEnvironment{quote}{\small}

% diseño de PIE DE PÁGINA
\usepackage{fancyhdr}
\pagestyle{fancy}
\fancyhf{}
\renewcommand{\headrulewidth}{0pt}
\rfoot{\thepage}
\fancypagestyle{plain}{\pagestyle{fancy}}

% DISEÑO DE LOS TÍTULOS de las partes del trabajo (capítulos y epígrafes o subcapítulos)
\usepackage{titlesec}
\usepackage{titletoc}
\titleformat{\chapter}[block]
{\large\bfseries\filcenter}
{\thechapter.}
{5pt}
{\MakeUppercase}
{}
\titlespacing{\chapter}{0pt}{0pt}{*3}
\titlecontents{chapter}
[0pt]                                               
{}
{\contentsmargin{0pt}\thecontentslabel.\enspace\uppercase}
{\contentsmargin{0pt}\uppercase}                        
{\titlerule*[.7pc]{.}\contentspage}                 

\titleformat{\section}
{\bfseries}
{\thesection.}
{5pt}
{}
\titlecontents{section}
[5pt]                                               
{}
{\contentsmargin{0pt}\thecontentslabel.\enspace}
{\contentsmargin{0pt}}
{\titlerule*[.7pc]{.}\contentspage}

\titleformat{\subsection}
{\normalsize\bfseries}
{\thesubsection.}
{5pt}
{}
\titlecontents{subsection}
[10pt]                                               
{}
{\contentsmargin{0pt}                          
	\thecontentslabel.\enspace}
{\contentsmargin{0pt}}                        
{\titlerule*[.7pc]{.}\contentspage}  

% NOTAS A PIE DE PÁGINA
\usepackage{chngcntr} % para numeración continua de las notas al pie
\counterwithout{footnote}{chapter}

% IMÁGENES
% DISEÑO DE FIGURAS. 
\usepackage{graphicx}
\graphicspath{{imagenes/}} % ruta a la carpeta de imágenes

%BIBLIOGRAFÍA 

% CONFIGURACIÓN PARA LA BIBLIOGRAFÍA IEEE
\usepackage[backend=biber, style=ieee, isbn=false,sortcites, maxbibnames=6, minbibnames=1]{biblatex} % Configuración para el estilo de citas de IEEE, recomendado para el área de ingeniería. "maxbibnames" indica que a partir de 6 autores trunque la lista en el primero (minbibnames) y añada "et al." tal y como se utiliza en el estilo IEEE.

% Añadimos las siguientes indicaciones para mejorar la adaptación del estilos en español
\DefineBibliographyStrings{spanish}{%
	andothers = {et\addabbrvspace al\adddot}
}
\DefineBibliographyStrings{spanish}{
	url = {\adddot\space[En línea]\adddot\space Disponible en:}
}
\DefineBibliographyStrings{spanish}{
	urlseen = {Acceso:}
}
\DefineBibliographyStrings{spanish}{
	pages = {pp\adddot},
	page = {p.\adddot}
}

\addbibresource{referencias.bib} % llama al archivo referencias.bib en el que deberá estar la bibliografía utilizada


%-------------
%	DOCUMENTO
%-------------

\begin{document}
\pagenumbering{roman} % Se utilizan cifras romanas en la numeración de las páginas previas al cuerpo del trabajo
	
%----------
%	PORTADA
%----------	
\begin{titlepage}
	%\begin{sffamily}
	\begin{center}
		\begin{Huge}
			<Título de la tesis>\\			
			\vspace{2cm}		
			<Nombre completo del doctorando/a>\\
		\end{Huge}
	    \vspace{4cm}	
	    \begin{Large}
            Tesis depositada en cumplimiento parcial de los requisitos para el grado de Doctor en\\
            \bigskip
            <Nombre del programa de doctorado>\\
		    \vspace{2cm}
		    Universidad Carlos III de Madrid\\
		    \vspace{3cm}
		    Director/a (es/as):\\
		    \bigskip
		    <Nombre completo Director/a 1>\\
		    <Nombre completo Director/a 2>\\
		    \bigskip
		    Tutor/a:\\
		    \bigskip
		    <Nombre completo Tutor/a>\\
		    \bigskip
		    <Mes de la defensa de tesis>\\
		 \end{Large}
		 
	\end{center}
	\end{titlepage}
	
%----------
%	LICENCIA CREATIVE COMMONS
%----------	
	% SI NUESTRO TRABAJO SE VA A PUBLICAR CON UNA LICENCIA CREATIVE COMMONS, INCLUIR ESTAS LÍNEAS. ES LA OPCIÓN RECOMENDADA.
	\thispagestyle{empty}
	\vspace*{\fill}
	\begin{center}
	Esta obra se encuentra sujeta a la licencia ``Creative Commons \textbf{Reconocimiento - No Comercial - Sin Obra Derivada}''.\\
	\medskip
	\includegraphics[width=4.2cm]{creativecommons.png} %incluimos el logotipo de Creative Commons
	\end{center}
	

%----------
%	DEDICATORIA
%----------	
% Si no quieres incluir una dedicatoria, elimina está sección.
\newpage
\thispagestyle{empty}
\begin{flushright}
Pon aquí tu dedicatoria.
\end{flushright}

%----------
%	AGRADECIMIENTOS
%----------	
% Si no quieres incluir agradecimientos, elimina está sección.
\chapter*{Agradecimientos}
\thispagestyle{empty}
Pon aquí tus agradecimientos. 

%----------
%	CONTENIDOS PUBLICADOS Y PRESENTADOS
%----------	
\chapter*{Contenidos publicados y presentados}
\thispagestyle{empty}
\begin{itemize}
    \item Relaciona una bibliografía de artículos u otras contribuciones de las que seas (co)autor que hayan sido incluidas como parte de la tesis y que
    \begin{itemize}
        \item hayan sido publicados (haz constar tu rol; las citas deben incluir DOIs o URLs del editor si están disponibles electrónicamente);  
        \item hayan sido presentados para su publicación (haz constar tu rol; las citas deben incluir  URLs si están disponibles electrónicamente);
        \item hayan sido presentados para la obtención de un título o de cualquier otra cualificación en la uc3m o en cualquier otra institución (haz constar el título o cualificación, la institución y el año de finalización).
    \end{itemize}
    \item En cada contribución 
    \begin{itemize}
        \item Declara si está incluida total o parcialmente en la tesis
        \item Indica el capítulo o capítulos en los que se incluye el material de dicha contribución
        \item Declaración opcional: “La inclusión de material de esta fuente en la tesis se especifica en una nota a pie de página en cada uno de los capítulos en los que se produce dicha inclusión”.
        \item Elige una de las siguientes declaraciones:
        \begin{itemize}
            \item El material de esta fuente incluido en la tesis no está señalado por medios tipográficos ni referencias.
            \item Todo material de esta fuente incluido en la tesis está señalado por medios tipográficos y una referencia explícita.
        \end{itemize}
        \end{itemize}
    \item En caso de que la tesis incluya cualquier material de terceros, incluidos trabajos de los que seas autor o coautor cuyos derechos hayas cedido, debes indicar que dicho material se usa con el correspondiente permiso. Por ejemplo: Fig.  2 reproducida con permiso del propietario de los derechos de autor, NombrePropietario
\end{itemize}
Importante: Si deseas hacer constar otros méritos de investigación, incluidas publicaciones y contribuciones no incluidas en la tesis, debes hacerlo en la página siguiente.

Si en la tesis no se incluye material de contribuciones publicadas o remitidas para su publicación, no incluyas esta página.

%----------
%	OTROS MÉRITOS DE INVESTIGACIÓN
%----------	
\chapter*{Otros méritos de investigación}
\thispagestyle{empty}
Puedes reseñar méritos de investigación, entre ellos artículos y contribuciones que, al no estar incluidos en la tesis, no están listados en la sección de CONTENIDOS PUBLICADOS Y PRESENTADOS.

Si no quieres reseñar otros méritos de investigación, no incluyas esta página. 



%----------
%	ÍNDICES
%----------	

%--
% Índice general
%-
\tableofcontents
\thispagestyle{fancy}

%--
% Índice de figuras. Si no se incluyen, comenta las líneas siguientes
%-
\listoffigures
\thispagestyle{fancy}

%--
% Índice de tablas. Si no se incluyen, comenta las líneas siguientes
%-
\listoftables
\thispagestyle{fancy}

%----------
%	MEMORIA
%----------	
\clearpage
\pagenumbering{arabic} % numeración con números arábigos para el resto de la memoria.	

\chapter{Introducción}

	% COMENZAR A ESCRIBIR la MEMORIA
	

%----------
%	BIBLIOGRAFÍA
%----------	

%\nocite{*} % Si quieres que aparezcan en la bibliografía todos los documentos que la componen (también los que no estén citados en el texto) descomenta está línea

\clearpage
\addcontentsline{toc}{chapter}{Bibliografía}
\setquotestyle[english]{british} % Cambiamos el tipo de cita porque en el estilo IEEE se usan las comillas inglesas.
\printbibliography



%----------
%	ANEXOS
%----------	

% Si tu trabajo incluye anexos, puedes descomentar las siguientes líneas
%\chapter* {Anexo x}
%\pagenumbering{gobble} % Las páginas de los anexos no se numeran



\end{document}